% =====================================================================
\chapter{Μοντέλο Απειλής και Σχετική Εργασία}
\label{ch:related}
% =====================================================================

\section{Ταξινόμηση επιθέσεων μέσω USB}
\label{sec:attack-taxonomy}

Οι \en{Nissim et~al.}~\cite{nissim2017usbattacks} καταγράφουν 29 διαφορετικές
επιθέσεις μέσω \en{USB} και τις ταξινομούν σε τέσσερις κατηγορίες, ενώ οι
\en{Tian et~al.}~\cite{tian2018sok} συστηματοποιούν τη γνώση ανά επίπεδο της
στοίβας \en{USB}. Για τους σκοπούς της εργασίας διακρίνουμε:

\begin{enumerate}
  \item \textbf{Επιθέσεις προσομοίωσης \en{HID}} (\en{BadUSB}, \en{Rubber Ducky},
        \en{O.MG}): η συσκευή δηλώνει πληκτρολόγιο και εγχέει εντολές.
  \item \textbf{Επιθέσεις μεταμφίεσης λειτουργιών}: σύνθετη συσκευή που κρύβει
        μια επικίνδυνη λειτουργία πίσω από μια αθώα (π.χ. «στικάκι» με κρυφό
        πληκτρολόγιο ή κάρτα δικτύου που ανακατευθύνει το \en{DNS}).
  \item \textbf{Επιθέσεις περιεχομένου}: κακόβουλα αρχεία ή σκόπιμα
        αλλοιωμένα συστήματα αρχείων που εκμεταλλεύονται σφάλματα του πυρήνα κατά
        την προσάρτηση.
  \item \textbf{Επιθέσεις κατά του οδηγού/πυρήνα}: αλλοιωμένοι περιγραφείς που
        εκμεταλλεύονται σφάλματα στον κώδικα απαρίθμησης.
  \item \textbf{Ηλεκτρικές επιθέσεις} (π.χ. \en{USB Killer}) και επιθέσεις
        \en{DMA} μέσω \en{Thunderbolt}.
\end{enumerate}

\section{Μοντέλο απειλής}
\label{sec:threat-model}

\subsection{Αντίπαλος}

Θεωρούμε αντίπαλο με \emph{σύντομη φυσική πρόσβαση} στον σταθμό εργασίας, ή
ικανότητα να πείσει τον χρήστη να συνδέσει μια συσκευή (π.χ. «χαμένο» στικάκι,
δωρεάν καλώδιο φόρτισης). Ο αντίπαλος ελέγχει πλήρως το υλικολογισμικό της
συσκευής: μπορεί να δηλώσει οποιουσδήποτε περιγραφείς, οποιοδήποτε όνομα
κατασκευαστή, να πληκτρολογήσει με οποιονδήποτε ρυθμό και να καθυστερήσει ή να
αλλάξει τη συμπεριφορά του.

\subsection{Υποθέσεις εμπιστοσύνης}

\begin{itemize}
  \item Ο υπολογιστής και ο πυρήνας του είναι αξιόπιστοι τη στιγμή που ξεκινά το
        \Probolos.
  \item Ο χρήστης απαντά στις ερωτήσεις με πληκτρολόγιο ή συνεδρία \en{SSH} που
        υπήρχε ήδη.
  \item Ο λογαριασμός \en{root} δεν έχει ήδη παραβιαστεί.
  \item Ο δίαυλος (ελεγκτές, διανομείς) είναι αξιόπιστος.
\end{itemize}

\subsection{Στόχοι ασφαλείας}

\begin{description}
  \item[Α1 -- Καμία σιωπηλή αποδοχή.] Καμία νέα συσκευή να μη λειτουργεί χωρίς
        ανθρώπινη απάντηση, απομνημονευμένη εμπιστοσύνη ή ρητή εξαίρεση.
  \item[Α2 -- Ενημερωμένη απόφαση.] Ο χρήστης να γνωρίζει τι δηλώνει η συσκευή
        και αν οι δηλώσεις της είναι αντιφατικές.
  \item[Α3 -- Αδυναμία αυτο-έγκρισης.] Μια συσκευή να μην μπορεί να απαντήσει
        στην ερώτηση που την αφορά.
  \item[Α4 -- Διαθεσιμότητα για τον νόμιμο χρήστη.] Το εργαλείο να μην μπορεί να
        κλειδώσει τον χρήστη έξω από το σύστημα.
\end{description}

\subsection{Εκτός πεδίου}

Ρητά εκτός πεδίου βρίσκονται: σφάλματα στον κώδικα \en{USB} του πυρήνα (τα
οποία ενεργοποιούνται πριν το \Probolos{} δει τη συσκευή)· κακόβουλα συστήματα
αρχείων· «υπομονετικές» συσκευές που επιτίθενται μέρες μετά την αποδοχή τους·
τέλεια αντίγραφα νόμιμων συσκευών που παραμένουν σιωπηλά κατά τον έλεγχο·
\en{Thunderbolt}/\en{DMA}· \en{USB Power Delivery}· \en{Bluetooth} και άλλες
ασύρματες συσκευές. Ο Πίνακας~\ref{tab:out-of-scope} εξηγεί τον λόγο για κάθε
περίπτωση.

\begin{table}[htbp]
\centering
\caption{Απειλές εκτός πεδίου και η αιτιολόγηση.}
\label{tab:out-of-scope}
\begin{tabular}{@{}p{4.2cm}p{9.6cm}@{}}
\toprule
\textbf{Απειλή} & \textbf{Γιατί δεν αντιμετωπίζεται}\\
\midrule
Σφάλματα \en{USB} πυρήνα & ο πυρήνας διαβάζει τους περιγραφείς πριν το εργαλείο δει τη συσκευή\\
Κακόβουλο σύστημα αρχείων & εκμεταλλεύεται τον πυρήνα κατά την προσάρτηση· το εργαλείο διαβάζει πίνακες διαμερισμάτων, όχι αρχεία\\
Υπομονετική συσκευή & περνά κάθε έλεγχο και επιτίθεται αργότερα· το ιστορικό ανιχνεύει μόνο αλλαγή περιγραφέων\\
Τέλειο αντίγραφο & δηλώνει ακριβώς ό,τι μια νόμιμη συσκευή· εντοπίζεται μόνο αν πληκτρολογήσει κατά την καραντίνα\\
\en{Thunderbolt}/\en{DMA} & διαφορετικό υποσύστημα· αρμοδιότητα του \en{IOMMU} και του \code{boltctl}\\
Ασύρματα & ένας προσαρμογέας ραδιοσυχνοτήτων περνά τους ελέγχους και δέχεται πληκτρολόγιο «από αέρος»\\
\bottomrule
\end{tabular}
\end{table}

\section{Σχετική εργασία}
\label{sec:related}

\subsection{Λύσεις λευκής λίστας: \en{USBGuard}}

Το \en{USBGuard}~\cite{usbguard} είναι το πιο διαδεδομένο εργαλείο ελέγχου
συσκευών \en{USB} στο \en{Linux}. Αξιοποιεί την ίδια διεπαφή εξουσιοδότησης και
εφαρμόζει πολιτικές κανόνων (π.χ. \code{allow with-interface equals \{ 08:*:* \}}).
Είναι ώριμο και ευέλικτο, αλλά η ασφάλειά του εξαρτάται πλήρως από την ποιότητα
της πολιτικής: δεν αξιολογεί αν οι ισχυρισμοί μιας άγνωστης συσκευής είναι
αντιφατικοί, δεν παρατηρεί συμπεριφορά και δεν εξετάζει τα μέσα αποθήκευσης.

\subsection{\en{GoodUSB}}

Το \en{GoodUSB}~\cite{tian2015goodusb} ζητά από τον χρήστη να δηλώσει τι
\emph{περιμένει} να είναι η συσκευή και περιορίζει τους οδηγούς σε αυτό.
Ύποπτες συσκευές ανακατευθύνονται σε εικονική μηχανή-«παγίδα» (\en{honeypot}).
Απαιτεί τροποποίηση του πυρήνα και εμπιστεύεται τη δήλωση του χρήστη, ο οποίος
συχνά δεν γνωρίζει τι είναι μια συσκευή.

\subsection{\en{USBFILTER}}

Το \en{USBFILTER}~\cite{tian2016usbfilter} λειτουργεί ως τείχος προστασίας
\emph{πακέτων} \en{USB} μέσα στον πυρήνα, με κανόνες τύπου \en{iptables} που
μπορούν να αναφέρονται και σε διεργασίες. Προσφέρει λεπτομερή έλεγχο, αλλά
απαιτεί τροποποιημένο πυρήνα και σύνταξη κανόνων από ειδικό.

\subsection{\en{Cinch}}

Το \en{Cinch}~\cite{angel2016cinch} απομονώνει τον δίαυλο \en{USB} σε ξεχωριστή
εικονική μηχανή και μεταφέρει την κίνηση μέσω μιας «πύλης» που επιβάλλει
πολιτικές και μπορεί να εφαρμόσει κρυπτογράφηση. Παρέχει ισχυρή απομόνωση (ακόμη
και από σφάλματα του πυρήνα), με κόστος πολυπλοκότητας και απαίτηση
εικονικοποίησης.

\subsection{\en{USBCheckIn}}

Το \en{USBCheckIn}~\cite{griscioli2016usbcheckin} απαιτεί από τον χρήστη να
αλληλεπιδράσει φυσικά με μια νέα συσκευή εισόδου (π.χ. να πληκτρολογήσει μια
τυχαία ακολουθία) πριν αυτή γίνει αποδεκτή, υλοποιημένο σε ενδιάμεσο υλικό.

\subsection{Ανίχνευση με βάση τον ρυθμό πληκτρολόγησης}

Το \en{USBlock}~\cite{neuner2018usblock} ανιχνεύει επιθέσεις έγχυσης από τον
χρονισμό των πληκτρολογήσεων, και το \en{USBESAFE}~\cite{kharraz2019usbesafe}
εκπαιδεύει ταξινομητή μηχανικής μάθησης πάνω σε χαρακτηριστικά της κίνησης
\en{USB}. Οι \en{Farhi et~al.}~\cite{farhi2019malboard} έδειξαν όμως με το
\en{Malboard} ότι μια συσκευή μπορεί να \emph{μιμηθεί} τον ρυθμό πληκτρολόγησης
συγκεκριμένου χρήστη, υπονομεύοντας ανιχνευτές που στηρίζονται μόνο στον χρονισμό.
Το εύρημα αυτό επηρέασε άμεσα τη σχεδίαση του σταδίου~3 (§\ref{sec:design-stage3}):
ο αποφασιστικός κανόνας δεν είναι ο ρυθμός, αλλά το γεγονός ότι η συσκευή
πληκτρολόγησε ενώ \emph{κανείς δεν την άγγιζε}.

\subsection{Συγκριτική επισκόπηση}

Ο Πίνακας~\ref{tab:related-comparison} συγκρίνει το \Probolos{} με τις
προαναφερθείσες προσεγγίσεις.

\begin{table}[htbp]
\centering
\small
\caption{Σύγκριση του \Probolos{} με σχετικές προσεγγίσεις
(\checkmark{} = υποστηρίζεται, \textbf{--} = όχι, $\sim$ = μερικώς).}
\label{tab:related-comparison}
\setlength{\tabcolsep}{4pt}
\begin{tabular}{@{}lccccccc@{}}
\toprule
 & \rotatebox{70}{\en{USBGuard}} & \rotatebox{70}{\en{GoodUSB}} & \rotatebox{70}{\en{USBFILTER}}
 & \rotatebox{70}{\en{Cinch}} & \rotatebox{70}{\en{USBCheckIn}} & \rotatebox{70}{\en{USBlock}} & \rotatebox{70}{\textbf{\Probolos}}\\
\midrule
Αμετάβλητος πυρήνας            & \checkmark & --         & --         & $\sim$     & \checkmark & \checkmark & \checkmark\\
Χωρίς πρόσθετο υλικό           & \checkmark & \checkmark & \checkmark & \checkmark & --         & \checkmark & \checkmark\\
Προεπιλεγμένη άρνηση           & \checkmark & \checkmark & $\sim$     & \checkmark & \checkmark & --         & \checkmark\\
Έλεγχος συνέπειας περιγραφέων  & --         & $\sim$     & --         & $\sim$     & --         & --         & \checkmark\\
Παρατήρηση συμπεριφοράς        & --         & $\sim$     & --         & --         & \checkmark & \checkmark & \checkmark\\
Έλεγχος μέσων αποθήκευσης      & --         & --         & --         & --         & --         & --         & \checkmark\\
Ανίχνευση μετάλλαξης           & $\sim$     & --         & --         & --         & --         & --         & \checkmark\\
Διαχωρισμός προνομίων          & $\sim$     & --         & --         & \checkmark & --         & --         & \checkmark\\
Προστασία από σφάλματα πυρήνα  & --         & --         & --         & \checkmark & $\sim$     & --         & --\\
\bottomrule
\end{tabular}
\end{table}

Η συνεισφορά του \Probolos{} δεν βρίσκεται σε κάποιον μεμονωμένο μηχανισμό, αλλά
στον \emph{συνδυασμό} τους σε ένα εργαλείο που εγκαθίσταται σε αμετάβλητη διανομή
\en{Linux}, με έμφαση στην ορθότητα των οριακών περιπτώσεων και στην ασφάλεια
έναντι κλειδώματος.
